% Source: https://artofproblemsolving.com/community/c3h2074734_simplification_how
% Stable identity: aops-topic-2074734-q01
Let \(x=a-c\) and \(y=b-d\). Then \(x+y=a+b-c-d=3\), and \(xy=ab-ad-bc+cd\). Adding the second and third equations gives \(4(ab-bc+cd-da)=9\), so \(xy=\dfrac{9}{4}\). Therefore
\[
(x-y)^2=(x+y)^2-4xy=9-9=0,
\]
so \(x=y=\dfrac{3}{2}\). (This is the step behind ``\((a+b-c-d)^2=4(ab-bc+cd-da)\) implies \((a-b-c+d)^2=0\)'': \(a-b-c+d=x-y\).) Hence
\[
11x^2+17y^2=28\cdot\dfrac{9}{4}=\boxed{63}.
\]
Such numbers exist: subtracting the second equation from the third gives \((a+c)(b+d)=\dfrac{1}{2}\), so for example \(a+c=1\), \(b+d=\dfrac{1}{2}\), \(a-c=b-d=\dfrac{3}{2}\) works.
